%% kapa: exam=1 round=32 year=2021 period=2 form=A subject=appraisal_law subject_ko=감정평가관계법규 seq=26 number=26
%% source: Q-Net 감정평가사 기출문제 (cst00309) artlSeq=5208397 file=extracted/1008/5208397/A형/2021년도 제32회 감정평가사 1차 2교시 A형.pdf
%% conversion: kordoc 4.24.51 -> tools/gen1.py (auto); review=proofread-sonnet; flags=none
\begin{question}{26}
건축법령상 건축허가에 관한 설명으로 옳은 것은? (단, 조례는 고려하지 않음)
\begin{choices}
  \item 21층 이상의 건축물을 특별시나 광역시에 건축하려면 국토교통부장관의 허가를 받아야 한다.
  \item 주거환경이나 교육환경 등 주변 환경을 보호하기 위하여 도지사가 필요하다고 인정하여 지정·공고한 구역에 건축하는 위락시설에 해당하는 건축물의 건축을 시장·군수가 허가하려면 도지사의 승인을 받아야 한다.
  \item 허가권자는 숙박시설에 해당하는 건축물의 건축을 허가하는 경우 해당 대지에 건축하려는 건축물의 용도·규모가 주거환경 등 주변 환경을 고려할 때 부적합하다고 인정되는 경우에는 건축위원회의 심의를 거치지 않고 건축허가를 하지 아니할 수 있다.
  \item 허가권자는 허가를 받은 자가 허가를 받은 날부터 4년 이내에 공사에 착수하지 아니한 경우라도 정당한 사유가 있다고 인정되면 2년의 범위에서 공사 기간을 연장할 수 있다.
  \item 분양을 목적으로 하는 공동주택의 건축허가를 받으려는 자는 대지의 소유권을 확보하지 않아도 된다.
\end{choices}
\end{question}
\answer{2}
